


Hypotheses, thresholds and kill criteria are in \texttt{docs/02\_preregistration/PREREGISTRATION\_V14.md} (committed before any T1, T2 or T3 run). Branch \texttt{final}; all algorithms are the frozen ones of \texttt{redesign-v1.1}; nothing was tuned. Labels: TRACE\_DRIVEN, FORMAL, MEASURED\_SOFTWARE, SIMULATED, ANALYTICAL, NOT\_RUN. No hardware data is used.

\subsection*{T1. Trace-driven replay on LoED (TRACE\_DRIVEN)}

Data: LoED (Zenodo record 4121430, CC BY 4.0; arXiv 2010.14211): 188 daily files, 8,055,823 frames kept after the CRC, counter and message-type filters, 8,503 device addresses seen. 2,975 devices have at least 300 frames and an admissible start position (5,363 have fewer than 300 frames; 165 have no segment of at least 41 frames). Split by hashed device address: 1,543 DEVELOPMENT, 1,432 HELD-OUT. The second dataset (Brno and Tampere authors, Zenodo 6228358, CC BY 4.0) holds one device (1,527 messages), so no split by device is possible; it is not used.

\subsubsection*{Link statistics (what the custom simulator assumed away)}

{\scriptsize\raggedright
\begin{longtable}{p{0.079\linewidth}p{0.079\linewidth}p{0.079\linewidth}p{0.079\linewidth}p{0.079\linewidth}p{0.079\linewidth}p{0.079\linewidth}p{0.079\linewidth}p{0.079\linewidth}p{0.079\linewidth}p{0.079\linewidth}p{0.079\linewidth}}
\toprule
Half & Devices & Loss below 10\% (L) & 10\% to 40\% (M) & 40\% or more (H) & Loss quantiles 5, 25, 50, 75, 95\% & Pooled loss & P(loss, next frame lost) & Ratio to pooled & Lag-1 autocorrelation & Mean burst (frames) & Bursts of 2 or more \\
\midrule
\endhead
DEVELOPMENT & 1,543 & 51 & 105 & 1,387 & 0.15, 0.67, 0.87, 0.93, 0.96 & 0.722 & 0.818 & 1.13 & 0.320 & 5.5 & 67\% \\
HELD-OUT & 1,432 & 45 & 90 & 1,297 & 0.15, 0.68, 0.87, 0.93, 0.96 & 0.702 & 0.821 & 1.17 & 0.391 & 5.6 & 67\% \\
\bottomrule
\end{longtable}}

Median RSSI -106 dBm and SNR -4.0 dB over all kept frames (5th percentile -115 dBm and -13.2 dB). Losses are positively correlated from frame to frame (a lost frame is followed by another lost frame more often than the pooled loss rate says), which the independent-loss simulator did not have. Most devices are weak links: the union over nine gateways still misses most of their frames, and 90\% of the HELD-OUT devices fall in group H. Group L holds 45 devices and group M 90, so their tests are the least precise.

\subsubsection*{Delivery of a Critical alert within 10 s on HELD-OUT devices (mean over devices of the mean over 300 start frames)}

{\scriptsize\raggedright
\begin{longtable}{p{0.158\linewidth}p{0.158\linewidth}p{0.158\linewidth}p{0.158\linewidth}p{0.158\linewidth}p{0.158\linewidth}}
\toprule
Policy & L (45 devices) & M (90 devices) & H (1,297 devices) & H, 95\% interval & Transmissions per alert, H \\
\midrule
\endhead
D1 & 99.8\% & 95.8\% & 49.3\% & 47.8\% to 50.8\% & 5.01 \\
quick retries & 99.7\% & 94.4\% & 42.7\% & 41.2\% to 44.1\% & 3.88 \\
severity table & 99.6\% & 91.5\% & 32.8\% & 31.6\% to 34.1\% & 3.26 \\
UA-SCT & 99.1\% & 87.4\% & 25.6\% & 24.5\% to 26.7\% & 1.84 \\
fixed 3 retries & 99.6\% & 91.5\% & 32.8\% & 31.6\% to 34.1\% & 3.26 \\
LoRaWAN-style confirmed uplink & 99.6\% & 91.7\% & 33.4\% & 32.1\% to 34.7\% & 2.76 \\
Nessa-style & 99.7\% & 94.4\% & 42.7\% & 41.2\% to 44.1\% & 4.06 \\
zero-spacing control (EXPLORATORY, not preregistered) & 99.9\% & 96.1\% & 54.0\% &  & 5.98 \\
\bottomrule
\end{longtable}}

\subsubsection*{Hypotheses (preregistered in V14)}

{\scriptsize\raggedright
\begin{longtable}{p{0.190\linewidth}p{0.190\linewidth}p{0.190\linewidth}p{0.190\linewidth}p{0.190\linewidth}}
\toprule
Hypothesis & Result & Interval & Raw p & Verdict (Holm, 7 tests) \\
\midrule
\endhead
H-T1a, group L: D1 not worse than the best baseline (quick) by more than 2 points & best minus D1 = -0.1 points & -0.2 to 0.0 points (upper limit at most 2) & below 0.001 & SUPPORTED \\
H-T1a, group M: D1 not worse than the best baseline (quick) by more than 2 points & best minus D1 = -1.4 points & -2.1 to -0.8 points (upper limit at most 2) & below 0.001 & SUPPORTED \\
H-T1a, group H: D1 not worse than the best baseline (quick) by more than 2 points & best minus D1 = -6.7 points & -6.8 to -6.5 points (upper limit at most 2) & below 0.001 & SUPPORTED \\
H-T1b: D1 beats quick retries by at least 2 points on group H & D1 minus quick = 6.7 points & 6.5 to 6.8 points & below 0.001 & SUPPORTED \\
H-T1c: the provisional alert of R3 stays within the R3 bounds on the traces (all three groups) & median earlier 1.87 s (L), 1.87 s (M), 1.87 s (H); extra airtime 37\%, 33\%, 31\% (upper limits 39\%, 35\%, 31\%); delivery difference 0.1, 1.8, 9.2 points & lower limits of the delivery difference 0.0, 1.2, 9.0 points & about 0.002 & SUPPORTED \\
\bottomrule
\end{longtable}}

\subsubsection*{How to read it (limits stated before the run and one found by the control)}

\begin{itemize}
\item \textbf{Delivery on the replay is decided by the number of attempts that finish before the deadline.} Attempt k sees the reception state of frame s + k, so the spacing of retries (seconds) and the spreading factor have no effect on reception, and nothing penalises an aggressive policy: there is no crowd, no collision and no duty sharing in a single-link replay. Quick retries and the Nessa-style policy have the same delivery on every device because both complete five attempts; the severity table and fixed 3 retries have the same delivery because both hold 5 s and complete four.
\item \textbf{The zero-spacing control} (retry at once, no hold, base spreading factor, up to 12 tries; added after the HELD-OUT run to separate retry timing from anything else, so it is EXPLORATORY and outside the Holm family) delivered 54.0\% on group H, above D1 (49.3\%). D1's gain over quick retries on the traces therefore comes from shorter retry spacing and more attempts, not from a mechanism that the traces can show; the traces cannot show the value of load-aware spacing, which is what D1's table is for.
\item The traces are uplink receptions, not ACK exchanges; downlink loss follows from the uplink loss rate (symmetric links); the devices are mostly static sensors, not vehicles; the union over nine gateways counts frames not heard by any gateway, so an idle gateway inflates loss; frames are seconds to hours apart, so loss at the scale of a retry is not observed.
\item The ns-3 result (D1 gain not reproduced, V12 addendum C) and this one agree on the point that matters for the claim: D1's advantage is not established beyond retry timing.
\end{itemize}

\subsection*{T2. Formal guarantees (FORMAL)}

\subsubsection*{T2.1 Security models (Tamarin 1.12.0, Maude 3.5.1, WSL Ubuntu 24.04, one thread)}

Models: \texttt{formal/crash\_alert\_auth.spthy} (final design M1, per-vehicle key kdf(master, vehicle), ACK bound to vehicle id and alert field; built with \texttt{-D=M0} for the earlier shared-key design, in which one cell key signs both alerts and ACKs) and \texttt{formal/crash\_alert\_fresh.spthy} (GPS-epoch freshness with a window of two epochs, an unbounded seen-set, sequence-in-second numbers s1 to s3, car reboot, base reboot, and the no-time fallback alert). Tool output: \texttt{formal/out/}. The models are symbolic with perfect cryptography; they say nothing about the implementation. The seen-set is unbounded (ring eviction is quantified in R1, not here), and the base counter check of the no-time alert is not modelled.

{\scriptsize\raggedright
\begin{longtable}{p{0.237\linewidth}p{0.237\linewidth}p{0.237\linewidth}p{0.237\linewidth}}
\toprule
Lemma & Model & Result & Expected (V14) \\
\midrule
\endhead
auth & M1 & verified & verified \\
L1 alert authenticity & M1 & verified & verified \\
L2 ACK authenticity per vehicle & M1 & verified & verified \\
L3 no cross-vehicle suppression & M1 & verified & verified \\
L1 on the shared-key design & M0 & falsified & falsified \\
L2 on the shared-key design & M0 & falsified & falsified \\
L3' cell-wide ACK suppression found on the shared-key design & M0 & falsified & falsified \\
executable (time-stamped alert accepted) & M1 & verified & verified \\
executable after a car reboot & M1 & verified & verified \\
L4 replay outside the window & M1 & verified & verified \\
L4b at most one acceptance while the seen-set is intact & M1 & verified & verified \\
L5 no double acceptance with a car reboot & M1 & verified & verified \\
L5 alert ids not reused after a car reboot & M1 & timeout & verified \\
L5 alert ids not reused after a car reboot, at most 3 ticks & M1 (3 ticks) & timeout & verified \\
L6 replay inside the window after a base reboot & M1 & falsified & falsified \\
L7 a no-time alert can be accepted epochs after it was sent & M1 & verified & verified \\
\bottomrule
\end{longtable}}

Reading: on M1 the tool proves that an accepted alert was sent by that vehicle (L1), that a car stops only on an ACK for its own accepted alert (L2), that a node compromised elsewhere cannot stop another vehicle (L3), that a time-stamped alert is accepted only in its own or the next epoch (L4), that no alert is accepted twice while the seen-set is intact (L4b) and that this holds across a car reboot. On M0 it finds the attack: with one cell key, a compromised node can sign an ACK for any alert and stop a vehicle that the base never heard (and can also sign alerts as another vehicle). L6 is refuted as expected: after a base reboot a replay inside the window is accepted again. L7 shows that a no-time alert can be accepted epochs after it was sent (no window for that mode).

\textbf{Not proved.} The lemma that an alert id (epoch and sequence number) sent before a car reboot is never reused after it did not finish: no result within 294 s (default search) and 421 s (at most three ticks); an induction variant of the earlier formulation timed out at 522 s. There is no proof and no counterexample, so H-T2sec as written (L1 to L5 proved) is NOT met; the part that is proved is stated above, and the argument for ids (boot takes at least one second, so the epoch differs) stays an argument, not a tool result. The numbers of the runs are in \texttt{formal/out/}.

\subsubsection*{T2.2 Timed guarantee (FORMAL by exhaustive enumeration; the tool is not UPPAAL)}

UPPAAL could not be used (academic licence and a registration step that this run cannot complete). The retry logic is a deterministic automaton over time on air, listen windows, holds and the duty gate, so the frozen policy code was enumerated exhaustively: D1 over all 4,096 failure histories of the 12 rounds before the alert, all loss patterns of the leading 8 attempts, and 0 to 35 s of airtime already used in the hour (\texttt{r3\_t2\_timed.py}). Random gaps take the end of their support (the worst case). Crash at 0, deadline 10 s, SF9, 1\% duty. Worst-case time from the crash to reception at the base when the initial k uplinks are lost (inf = not delivered before the deadline):

{\scriptsize\raggedright
\begin{longtable}{p{0.119\linewidth}p{0.119\linewidth}p{0.119\linewidth}p{0.119\linewidth}p{0.119\linewidth}p{0.119\linewidth}p{0.119\linewidth}p{0.119\linewidth}}
\toprule
Policy & k = 0 & 1 & 2 & 3 & 4 & 5 & 6 \\
\midrule
\endhead
D1 & 2.38 s & 3.32 s & 4.71 s & 6.31 s & 7.46 s & 8.39 s & 9.79 s \\
quick retries & 2.38 s & 3.62 s & 5.31 s & 7.21 s & 9.11 s & inf & inf \\
severity table & 7.38 s & 8.32 s & 9.26 s & inf & inf & inf & inf \\
UA-SCT & 2.38 s & 8.32 s & inf & inf & inf & inf & inf \\
fixed 3 retries & 7.38 s & 8.32 s & 9.26 s & inf & inf & inf & inf \\
LoRaWAN-style confirmed uplink & 2.38 s & 6.32 s & inf & inf & inf & inf & inf \\
Nessa-style & 2.38 s & 3.45 s & 5.39 s & 7.44 s & inf & inf & inf \\
\bottomrule
\end{longtable}}

H-T2: D1's worst case is not longer than quick retries' for every k from 0 to 6 (holds); no run of 134,144 exceeded the 36 s hourly allowance (0 breaches). Quick retries cannot deliver within 10 s after 5 or 6 initial losses; D1 can (9.79 s at k = 6), because it retries at once and keeps the base spreading factor in the worst case. This is a statement about the frozen policy code in a single-link model with the stated timing, not about a radio; the check agrees with the trace result that D1's advantage is more attempts inside the deadline.

\subsection*{T3. Two-stage alert on real crash data (MEASURED\_SOFTWARE)}

Data: the de-duplicated VZCrash test split, stored per-event decisions of \texttt{s6\_early} (seed 42; thresholds chosen on validation; no refit; the split was already scored once by \texttt{s6\_early} and used again by R3), and the NHTSA pulses of \texttt{nhtsa.py}.

{\scriptsize\raggedright
\begin{longtable}{p{0.158\linewidth}p{0.158\linewidth}p{0.158\linewidth}p{0.158\linewidth}p{0.158\linewidth}p{0.158\linewidth}}
\toprule
Target & Triggered events & Crash events that alarm in both stages & Median provisional decision & Median earlier than the full decision (95\% interval) & Retraction share (95\% interval) \\
\midrule
\endhead
1\% false alarms & 5,750 & 2,397 & 50 ms & 1750 ms (1750 to 1750) & 10.8\% (9.7\% to 12.0\%) \\
5\% false alarms & 5,750 & 3,809 & 100 ms & 1700 ms (1700 to 1700) & 1.2\% (0.9\% to 1.5\%) \\
\bottomrule
\end{longtable}}

At the 1\% target 67\% of those crash events are decided at the 50 ms checkpoint, 12\% at 100 ms, 6\% at 200 ms, 3\% at 400 ms and 5\% at 800 ms. H-T3a (SUPPORTED): the median provisional decision is 1750 ms earlier than the full decision (lower limit 1750 ms, bound 800 ms). H-T3b (SUPPORTED): the retraction share is 10.8\% with an upper limit of 12.0\%, bound 20\%.

\subsubsection*{Severity class of the provisional window against the full window}

VZCrash: 2,647 crash events with a provisional alarm; class agreement 99.8\% (99.7\% to 100.0\%). This carries almost no information: VZCrash has no severity labels and the computed full-window class is 0 for nearly every event (counts of classes 0 to 3: \{"0": 2643, "1": 4, "2": 0, "3": 0\}).

NHTSA pulses (sensor model at 16 g and 100 Hz; class from \texttt{severity.py}), share of pulses whose class from the window truncated at the given length equals the class from the full window:

{\scriptsize\raggedright
\begin{longtable}{p{0.119\linewidth}p{0.119\linewidth}p{0.119\linewidth}p{0.119\linewidth}p{0.119\linewidth}p{0.119\linewidth}p{0.119\linewidth}p{0.119\linewidth}}
\toprule
Pulses & n & 50 ms & 100 ms & 200 ms & 400 ms & 800 ms & Weighted by the VZCrash decision times \\
\midrule
\endhead
frontal barrier & 116 & 3\% & 36\% & 98\% & 100\% & 100\% & 19\% \\
side impactor & 60 & 52\% & 97\% & 100\% & 100\% & 100\% & 60\% \\
\bottomrule
\end{longtable}}

The two-stage alert is early because most provisional decisions are made at 50 ms, and at that length the truncated window does not yet carry the severity class of a frontal crash (it agrees for about 3\% of the frontal pulses); the 14-byte flash of R3 does not carry a class, so the design does not depend on it, but a severity-carrying provisional alert would have to wait for 200 ms. Time to decision on the NHTSA pulses is NOT\_RUN: no early-exit detector is trained on them, and applying the VZCrash-trained forest would be an unplanned transfer test.

H-T3c (SUPPORTED): the end-to-end earlier arrival (median at least 0.8 s) holds in 2 of 3 link models: the custom simulator (\texttt{ar.r3}, SIMULATED), and the trace-driven replay (T1c, TRACE\_DRIVEN, median 1.87 s); ns-3 is NOT\_RUN (the ns-3 application has no two-stage alert) and counts as not holding.

\subsection*{T4. Post-crash survivability budget (ANALYTICAL)}

Inputs read in opened documents (2026-11-15): SX1276/77/78/79 datasheet Rev. 4 (Semtech, March 2015; copy hosted by Adafruit): transmit current 120 mA at +20 dBm and 87 mA at +17 dBm (PA\_BOOST), 29 mA at +13 dBm and 20 mA at +7 dBm (RFO), LoRa receive 11.5 mA at 125 kHz; ESP32 Series Datasheet v5.3 (Espressif), Table 4-2: CPU on with the radio off, 160 MHz, 27 to 44 mA (44 mA used), deep sleep 10 uA; u-blox NEO-6 data sheet GPS.G6-HW-09005-E, Table 11: NEO-6M tracking 39 mA, acquisition 47 mA at 3.0 V. ASSUMPTIONS, not read in an opened document: 3.3 V at the three parts; converter efficiency 0.85; a capacitor usable from 5.0 V down to 3.0 V; no ESR loss; GPS already tracking; the node awake from the crash to the end of the schedule. No hypothesis is tested.

Energy drawn from the store for one Critical alert schedule at +17 dBm (all uplinks lost up to the 10 s deadline; and delivery after k lost uplinks), and the smallest capacitor that supplies the all-lost schedule:

{\scriptsize\raggedright
\begin{longtable}{p{0.119\linewidth}p{0.119\linewidth}p{0.119\linewidth}p{0.119\linewidth}p{0.119\linewidth}p{0.119\linewidth}p{0.119\linewidth}p{0.119\linewidth}}
\toprule
Policy & All lost: energy & Transmissions & Smallest capacitor & k = 0 & k = 2 & k = 4 & k = 6 \\
\midrule
\endhead
D1 & 4.97 J & 8 & 0.62 F & 0.95 J & 1.95 J & 2.95 J & 3.95 J \\
quick retries & 4.70 J & 5 & 0.59 F & 0.95 J & 2.44 J & 4.38 J & 4.70 J (not delivered) \\
severity table & 4.08 J & 4 & 0.51 F & 2.56 J & 3.56 J & 4.08 J (not delivered) & 4.08 J (not delivered) \\
UA-SCT & 3.62 J & 2 & 0.45 F & 0.95 J & 3.62 J (not delivered) & 3.62 J (not delivered) & 3.62 J (not delivered) \\
fixed 3 retries & 4.08 J & 4 & 0.51 F & 2.56 J & 3.56 J & 4.08 J (not delivered) & 4.08 J (not delivered) \\
LoRaWAN-style confirmed uplink & 3.90 J & 3 & 0.49 F & 0.95 J & 3.90 J (not delivered) & 3.90 J (not delivered) & 3.90 J (not delivered) \\
Nessa-style & 4.55 J & 5 & 0.57 F & 0.95 J & 2.31 J & 4.55 J (not delivered) & 4.55 J (not delivered) \\
\bottomrule
\end{longtable}}

Usable energy of the stated capacitors: 0.22 F: 1.76 J, 0.47 F: 3.76 J, 1.0 F: 8.00 J, 2.2 F: 17.60 J. Most of the energy is the ESP32 and the GPS staying awake for the 10 s window (about 3.2 J), not the radio: the policies differ by about 1 J. At +17 dBm a 1 F capacitor supplies every policy's full schedule; a 0.47 F capacitor supplies none of the all-lost schedules but supplies a delivery after up to six losses for D1 (3.95 J). At +20 dBm the all-lost schedules need 3.8 J to 5.6 J. This depends on the assumed converter and voltage window and on keeping the GPS powered.

\subsection*{Deviations from the preregistration and slips (also in DECISIONS.md)}

\begin{itemize}
\item T2.2 used exhaustive enumeration of the frozen code instead of UPPAAL (not obtainable); TChecker was not built (time).
\item The L5 id-uniqueness lemma was reformulated once after the earlier formulation timed out (the formulation in V14 was the statement that alert ids are unique across reboots); neither formulation finished.
\item A zero-spacing control policy was added to T1 after the HELD-OUT run and is labelled EXPLORATORY; the preregistered numbers were recomputed after its addition and are unchanged (all 167 stored keys equal).
\item \texttt{replay\_alert} gained two optional arguments (initial airtime used, a transmission trace) after the T1 run, for T2.2 and T4; the T1 results were recomputed and are unchanged.
\item The T1 population is dominated by weak links (90\% of HELD-OUT devices in group H); the group limits were fixed before the loss rates were computed and were not changed.
\item T3 reuses the stored seed-42 decisions on the de-duplicated test split, which \texttt{s6\_early} and R3 had already used; no model was refitted.
\end{itemize}

